\chapter{Motor de reglas: ZINA}

Reglas sencillas, narrativas y de supervivencia para NeoBilbao.

\section{Las estadísticas emocionales}

En NeoBilbao no existen los atributos físicos tradicionales. Las acciones de los personajes se resuelven utilizando sus emociones actuales como modificadores (de $-1$ a $+2$). Cómo te sientes determina cómo reaccionas ante el peligro:

\begin{itemize}
  \item \lead{Enfadado:} rebeldía, fuerza bruta, confrontación directa, romper obstáculos o enfrentarse cara a cara a una amenaza.
  \item \lead{Asustado:} supervivencia pura, huida, esquivar, reaccionar por instinto o buscar refugio a ciegas.
  \item \lead{Culpable:} el peso de la mentira, la resistencia mental, hackear sistemas corporativos o soportar los ataques psicológicos de Sugaar.
  \item \lead{Afligido:} empatía, conectar con otros, descifrar verdades ocultas, mantener la humanidad o consolar a los aliados.
  \item \lead{Pacífico:} calma calculada, análisis sereno del entorno, infiltración silenciosa y sigilo.
\end{itemize}

\begin{CajaRegla}{Nota de campaña: la abstinencia del Etxea}
Al romperse el Velo, el estado \emph{Alegre} desaparece permanentemente de la hoja del personaje. Ya no hay espacio para la euforia artificial.
\end{CajaRegla}

\section{La mecánica de resolución y el uso de etiquetas}

Cuando un personaje realiza una acción arriesgada descrita en la ficción, tira 2d6 + el modificador de la estadística emocional que corresponda:

\begin{itemize}
  \item \lead{10 o más (éxito total):} consigues exactamente lo que querías sin complicaciones, aplicando las ventajas de las etiquetas de tu equipo o entorno.
  \item \lead{7 a 9 (éxito con coste):} logras tu objetivo, pero el Director (MC) introduce una complicación. El Director puede gastar o activar las etiquetas del escenario o del equipo (por ejemplo, haciendo que un arma con etiqueta \etiqueta{Consumible} se agote, o que el entorno con etiqueta \etiqueta{Ácido} muerda el equipo del personaje).
  \item \lead{6 o menos (fallo):} las cosas salen rematadamente mal. Pierdes la iniciativa y el Director ejecuta un movimiento duro.
\end{itemize}

\section{Zama: estrés, carga y descompresión}

Los personajes no tienen barras de vida tradicionales. En su lugar, acumulan Zama (carga/estrés) en una escala del 0 al 6.

\lead{Cómo se gana.} Al recibir ataques físicos, presenciar horrores (como la sala de tanques en la Sesión~4), o fallar tiradas clave ante la presión de Sugaar.

\lead{Cómo se alivia (descomprimir).} Cuando el grupo encuentra un refugio seguro en los Zócalos (un conducto abandonado o un escondite de la resistencia) y dedica un momento a hablar, compartir una ración o calmar los nervios, pueden reducir 1 punto de Zama describiendo qué les consuela o qué recuerdo humano les ancla a la realidad.

\lead{El límite del colapso.} Si un personaje llega a 6 de Zama, sufre un colapso. El Director le impondrá una consecuencia drástica en la ficción (quedar fuera de combate, huir presa del pánico absoluto, desmayarse o obedecer ciegamente un mandato del sistema por agotamiento mental).

\section{Auzolan: reserva colaborativa}

Representa la solidaridad obrera heredada del Auzolan industrial. El grupo comparte una reserva común de puntos de Auzolan (empiezan con 1 punto por jugador al inicio de la sesión).

\begin{itemize}
  \item \lead{Ayudar a un compañero:} gastas 1 punto de Auzolan para dar un $+1$ a la tirada de otro jugador, describiendo cómo te pones en peligro o le animas.
  \item \lead{Declarar un hecho:} gastas 1 punto de Auzolan para alterar la narrativa menor («Conozco este sector porque trabajaba aquí», «Tengo un soplete escondido en esta tubería»).
\end{itemize}

\section{Movimientos básicos}

Disponibles para todos los PJs.

\begin{itemize}
  \item \lead{Atacar / Forcejear (usando Enfadado o Asustado).} Cuando te enfrentas a un peligro físico directo o agredes a un enemigo, tira con tu emoción dominante. Con 10+, infliges daño o superas el obstáculo sin recibir represalias. Con 7--9, logras tu objetivo o dañas al rival, pero el Director impone un coste (recibes daño/estrés, tu equipo sufre desgaste o el enemigo contraataca). Con 6 o menos, fallas y el Director ejecuta un movimiento duro.
  \item \lead{Investigar / Analizar (usando Pacífico o Culpable).} Cuando examinas con atención un lugar, un cadáver o un terminal en busca de pistas, tira con tu estado. Con 10+, haz 2 preguntas al Director sobre la escena. Con 7--9, haz solo 1 pregunta, pero el Director te revelará un detalle a cambio de una complicación o estrés.
  \item \lead{Manipular / Convencer (usando Afligido o Enfadado).} Cuando intentas que un NPC (un obrero anestesiado, un Bitartekoa) haga lo que quieres apelando a sus sentimientos o a la rabia, tira. Con 10+, accede a tus demandas. Con 7--9, accede, pero exige algo a cambio o duda en el último momento.
  \item \lead{Infiltrarse / Ocultarse (usando Pacífico o Asustado).} Cuando intentas moverte sin que te detecten los Tártalos o las cámaras, tira. Con 10+, pasas inadvertido en las sombras. Con 7--9, el sigilo se tambalea; haces ruido o dejas rastro, obligando a moverte rápido.
\end{itemize}

\section{El sistema de etiquetas (tags)}

En el motor «Zina», las etiquetas son palabras clave breves que definen las propiedades mecánicas y narrativas del equipo, el entorno o la condición de los personajes. Permiten dinamizar la ficción sin recurrir a cálculos complejos.

\newcommand{\tagtable}[1]{%
  \rowcolors{1}{white}{BoxGray}%
  \begin{tabularx}{\linewidth}{@{\hspace{3pt}}>{\raggedright\arraybackslash}p{.30\linewidth}>{\raggedright\arraybackslash}X@{}}
  #1
  \end{tabularx}\par\medskip}

\subsection{A. Etiquetas de equipo y armas (los Zócalos)}

Las herramientas improvisadas y el armamento de los obreros o de la seguridad corporativa se rigen por estas propiedades:

\tagtable{%
\etiqueta{Contundente} & Ideal para romper blindajes o deformar chapa, pero ruidoso. \\
\etiqueta{Desgarrador} & Causa heridas graves y desangramiento; aumenta el riesgo de elevar la Zama. \\
\etiqueta{A distancia corta} & Solo es efectivo en espacios reducidos o cuerpo a cuerpo. \\
\etiqueta{Pesado} & Requiere esfuerzo físico o penaliza las tiradas de sigilo (Pacífico o Asustado). \\
\etiqueta{Consumible} & El objeto tiene un uso limitado y se agota tras la tirada (como un soplete o un inyector). \\
\etiqueta{Clandestino} & Fácil de ocultar de los escáneres de los drones Tártalos. \\
\etiqueta{Preciso / A distancia} & Permite golpear objetivos lejanos o vulnerables (como los sensores ópticos de los drones), pero es inútil en forcejeos cuerpo a cuerpo. \\
\etiqueta{Improvisado} & Creado con chatarra de la fundición; se rompe o degrada automáticamente tras sacar un éxito con coste (7--9) o un fallo. \\}

\subsection{B. Etiquetas de entorno (la Ría de Acero y la Cúpula)}

El Director (MC) puede invocar estas etiquetas del escenario para justificar complicaciones en los éxitos con coste o los movimientos duros:

\tagtable{%
\etiqueta{Ácido} & El entorno corroe la materia orgánica y el metal; daña el equipo o la salud si te detienes. \\
\etiqueta{Oscuro / Niebla} & Ofrece cobertura natural para el sigilo, pero limita la visión y facilita las emboscadas. \\
\etiqueta{Vibración constante} & El ruido de la maquinaria oculta pasos y sonidos, pero dificulta las comunicaciones por radio. \\
\etiqueta{Alta temperatura} & El calor extremo fatiga a los personajes y acelera la ganancia de Zama. \\
\etiqueta{Vigilado} & Las cámaras o los sensores de ADUR-OS están activos en el sector. \\
\etiqueta{Inestable / Frágil} & Pasarelas oxidadas o estructuras a punto de ceder; un paso en falso activa una caída o un derrumbe. \\
\etiqueta{Electrificado} & Cables pelados o fugas de plasma energético; cruzar sin protección expone a daños eléctricos severos. \\}

\subsection{C. Etiquetas de condición (estados temporales)}

Como consecuencia de movimientos duros o ataques especiales de los enemigos, los personajes pueden recibir etiquetas de estado que limitan sus acciones hasta que sean curadas o superadas:

\tagtable{%
\etiqueta{Atemorizado} & Impide usar estadísticas de confianza o calma; penaliza las tiradas de resolución. \\
\etiqueta{Intoxicado} & Los pulmones sufren por el gas o el sirimiri ácido; penaliza las acciones físicas. \\
\etiqueta{Desarmado} & Has perdido tu herramienta o arma principal en el fragor de la huida. \\
\etiqueta{Aislado} & Separado del grupo por una compuerta o un derrumbe; no puedes recibir apoyo de Auzolan. \\
\etiqueta{Aturdido} & La mente zumba por una descarga o un eco de Sugaar; pierdes la iniciativa en el próximo intercambio de acciones. \\
\etiqueta{Herido / Cojeando} & Tu movilidad está castigada; no puedes huir con éxito sin ayuda o sin sacrificar equipo. \\}
